\documentclass[11pt]{article}
\usepackage{amsmath,amssymb,amsthm}
\usepackage{mathtools}
\usepackage{graphicx}
\usepackage{hyperref}
\usepackage{geometry}
\geometry{a4paper, margin=1in}

\newtheorem{theorem}{Theorem}[section]
\newtheorem{lemma}[theorem]{Lemma}
\newtheorem{proposition}[theorem]{Proposition}
\newtheorem{corollary}[theorem]{Corollary}
\newtheorem{definition}[theorem]{Definition}
\newtheorem{remark}[theorem]{Remark}
\newtheorem{example}[theorem]{Example}

\newcommand{\C}{\mathbb{C}}
\newcommand{\R}{\mathbb{R}}
\newcommand{\N}{\mathbb{N}}
\newcommand{\Z}{\mathbb{Z}}
\newcommand{\iu}{\mathrm{i}} % imaginary unit
\newcommand{\e}{\mathrm{e}}
\DeclareMathOperator{\ReOp}{Re}
\DeclareMathOperator{\ImOp}{Im}
\DeclareMathOperator{\argmax}{arg\,max}
\DeclareMathOperator{\sgn}{sgn}

\title{The Iota Function:\\ A Geometric Encoding of the Dirichlet Eta Series}
\author{Generated by AI Assistant}
\date{\today}

\begin{document}

\maketitle

\begin{abstract}
We introduce the \emph{iota function}, a mapping that encodes a complex variable $s = \sigma + \iu t$ into a sequence of points in $\R^3$ (discrete case) or a continuous space curve (continuous case). The construction preserves the term-by-term structure of the Dirichlet eta series $\eta(s) = \sum_{n=1}^\infty (-1)^{n-1} n^{-s}$ in an invertible manner, thereby providing a faithful geometric representation of the series. We establish the injectivity of both versions, derive explicit recovery formulas for $s$, and discuss the relationship between the iota spiral and the Riemann zeta function. In particular, we translate the Nyman--Beurling--B\'aez-Duarte greedy condition for the Riemann Hypothesis into a geometric constraint on the continuous iota curve. Applications to complex analysis, dynamical systems, and the geometry of special functions are outlined.
\end{abstract}

\tableofcontents

\section{Introduction}

The Dirichlet eta function,
\begin{equation}
  \eta(s) = \sum_{n=1}^{\infty} (-1)^{n-1} n^{-s}, \qquad \ReOp(s) > 0,
\end{equation}
is an alternating variant of the Riemann zeta function $\zeta(s)$, related by the identity
$\eta(s) = (1 - 2^{1-s})\,\zeta(s)$.  The non‑trivial zeros of $\eta(s)$ coincide with those of $\zeta(s)$, so any result concerning the zeros of $\eta(s)$ is directly relevant to the Riemann Hypothesis (RH).

Each term of the eta series can be viewed as a vector in the complex plane:
$v_n(s) = (-1)^{n-1} n^{-s} = n^{-\sigma} \e^{-\iu t \ln n}$ (with an unwrapped phase).  The sum $\eta(s)$ collapses the entire sequence $\{v_n\}$ into a single complex number, discarding the term‑by‑term information.  The \emph{iota function}, introduced in this note, restores this lost information by lifting the sequence of terms to a three‑dimensional space curve.  The construction is both discrete (evaluated at integers) and continuous (extended to a real parameter).

The iota function is not merely a visualisation tool; it is \emph{injective}, meaning that different values of $s$ produce geometrically distinct objects.  Consequently, any property of $\eta(s)$—including the Riemann Hypothesis—can be reformulated as a geometric condition on the iota spiral.  We develop this connection and discuss its potential implications across several branches of mathematics.

\section{Definition of the Iota Function}

\subsection{Discrete Iota Function}

\begin{definition}[Discrete Iota]\label{def:iota_discrete}
  For a complex number $s = \sigma + \iu t$ with $\sigma > 0$, define for each integer $n \ge 1$ the quantities
  \begin{align}
    r_n &= n^{-\sigma}, \label{eq:radius}\\
    \Phi_n &= \pi(n-1) - t \ln n \qquad \text{(unwrapped phase)}. \label{eq:phase}
  \end{align}
  The \emph{discrete iota function} $\iota(s)$ is the sequence of points $P_n \in \R^3$ given by
  \begin{equation}
    P_n = \bigl( r_n \cos\Phi_n,\; r_n \sin\Phi_n,\; n \bigr), \qquad n = 1,2,3,\dots \label{eq:iota_disc}
  \end{equation}
\end{definition}

The first two coordinates of $P_n$ are exactly the real and imaginary parts of the $n$-th term of the eta series, while the third coordinate is simply the index $n$.  The points lie on a discrete spiral that winds around the vertical $z$-axis.  Because $\Phi_n$ is the \emph{unwrapped} total accumulated angle (not taken modulo $2\pi$), the sequence carries strictly more information than the projected complex numbers $v_n(s)$.

\subsection{Continuous Iota Function}

To obtain a smooth curve that interpolates the discrete points, we replace the integer index $n$ by a continuous real parameter $\tau \ge 1$.

\begin{definition}[Continuous Iota]\label{def:iota_continuous}
  For $s = \sigma + \iu t$ with $\sigma > 0$, the \emph{continuous iota curve} $\iota_c(s) : [1,\infty) \to \R^3$ is defined by
  \begin{equation}
    \iota_c(\tau) = \Bigl(
      \tau^{-\sigma} \cos\bigl(\pi(\tau-1) - t \ln \tau\bigr),\;
      \tau^{-\sigma} \sin\bigl(\pi(\tau-1) - t \ln \tau\bigr),\;
      \tau
    \Bigr). \label{eq:iota_cont}
  \end{equation}
\end{definition}

By construction, $\iota_c(n) = P_n$ for every integer $n \ge 1$, so the continuous curve interpolates the discrete iota points exactly.  The curve is $C^\infty$ for $\tau > 0$ and provides a faithful geometric encoding of $s$.

\section{Injectivity and Recovery of $s$}

Both versions of the iota function are injective; the complex number $s$ can be uniquely reconstructed from the geometric object.

\begin{theorem}[Injectivity of Discrete Iota]\label{thm:inj_disc}
  The map $s \mapsto \iota(s)$ is injective.  Given the sequence $\{P_n\}_{n=1}^\infty$, the parameters $\sigma$ and $t$ can be recovered explicitly:
  \begin{align}
    \sigma &= -\frac{\ln r_2}{\ln 2}, \label{eq:rec_sigma}\\
    t &= \frac{\pi(n-m) - (\Phi_n - \Phi_m)}{\ln(n/m)}, \quad \text{for any } n \neq m. \label{eq:rec_t}
  \end{align}
\end{theorem}
\begin{proof}
  From $P_2 = (r_2\cos\Phi_2, r_2\sin\Phi_2, 2)$ we read off $r_2 = 2^{-\sigma}$, which immediately yields \eqref{eq:rec_sigma}.  With $\sigma$ known, the unwrapped phase $\Phi_n$ is determined (up to a global multiple of $2\pi$ that cancels in the difference).  For any distinct indices $n,m$, the difference of phases is
  \[
    \Phi_n - \Phi_m = \pi(n-m) - t \ln\frac{n}{m},
  \]
  from which \eqref{eq:rec_t} follows.  The formula is valid for any pair because the unwrapped phases are directly accessible from the continuous tracking of the angles of the points $P_n$.
\end{proof}

\begin{theorem}[Injectivity of Continuous Iota]\label{thm:inj_cont}
  The map $s \mapsto \iota_c(s)$ is injective.  Given the curve $\gamma(\tau) = \iota_c(\tau)$, one can recover $\sigma$ and $t$ from any two parameter values $\tau_1 \neq \tau_2$ using the radial and angular data.
\end{theorem}
\begin{proof}
  Let $r(\tau) = \tau^{-\sigma}$ be the distance from the $z$-axis.  Then
  \[
    \frac{r(\tau_1)}{r(\tau_2)} = \left(\frac{\tau_1}{\tau_2}\right)^{-\sigma} \implies \sigma = -\frac{\ln(r(\tau_1)/r(\tau_2))}{\ln(\tau_1/\tau_2)}.
  \]
  The unwrapped phase difference is
  \[
    \phi(\tau_1) - \phi(\tau_2) = \pi(\tau_1-\tau_2) - t \ln\frac{\tau_1}{\tau_2},
  \]
  whence $t$ is obtained.  The unwrapped phase can be tracked continuously along the curve because the curve is smooth and the parameterisation is explicit.
\end{proof}

\section{Properties of the Iota Spiral}

\subsection{Decay and Asymptotic Behaviour}
For $\sigma > 0$, the radius $r(\tau) = \tau^{-\sigma}$ decreases monotonically, so the spiral converges to the $z$-axis as $\tau \to \infty$.  The speed of the curve is
\begin{equation}
  \|\iota_c'(\tau)\| = \sqrt{1 + \tau^{-2\sigma}\Bigl(\sigma^2\tau^{-2} + \bigl(\pi - \tfrac{t}{\tau}\bigr)^2\Bigr)}.
\end{equation}
As $\tau \to \infty$, the speed tends to $1$, reflecting the fact that the curve asymptotically becomes a vertical line with vanishing horizontal oscillations.

\subsection{Curvature and Torsion}
The differential geometry of $\iota_c$ can be analysed via the Frenet--Serret apparatus.  Let $\kappa(\tau)$ be the curvature and $\tau_{\text{tor}}(\tau)$ the torsion.  Explicit formulas show that both quantities depend non‑trivially on $\sigma$, $t$, and $\ln\tau$.  In particular, the torsion does not vanish identically, indicating that the curve is genuinely three‑dimensional and not confined to a plane.

\subsection{Length and Convergence}
The total arc length of $\iota_c$ from $\tau = 1$ to $\infty$ is finite if and only if $\sigma > 1$.  For $\sigma = 1/2$ (the critical line), the horizontal oscillations contribute an infinite arc length, yet the projection onto the $xy$-plane converges conditionally.  This mirrors the conditional convergence of the Dirichlet eta series on the critical line.

\section{Relation to the Dirichlet Eta and Riemann Zeta Functions}

The iota function is intimately tied to the Dirichlet eta function via summation formulas.

\subsection{Abel--Plana Representation}
Applying the Abel--Plana formula to $f(\tau) = \e^{\iu\pi(\tau-1)} \tau^{-s}$ yields
\begin{equation}
  \eta(s) = \frac{1}{2} + \int_1^\infty \e^{\iu\pi(\tau-1)} \tau^{-s}\,d\tau + \mathcal{R}(s),
  \label{eq:abel_plana}
\end{equation}
where $\mathcal{R}(s)$ is an absolutely convergent remainder integral.  The integral term is exactly the continuous accumulation of the horizontal projections of the iota curve:
\[
  \int_1^\infty \bigl( x(\tau) + \iu y(\tau) \bigr) \,d\tau,
\]
with $(x(\tau), y(\tau), z(\tau)) = \iota_c(\tau)$.  Thus $\eta(s)$ is, up to a bounded correction, the projected area (or complex displacement) of the continuous iota spiral.

\subsection{The Greedy Condition and the Riemann Hypothesis}
A well‑known equivalent formulation of the Riemann Hypothesis (for the eta function) is the \emph{greedy condition}: if $\eta(s) = 0$ and $\ReOp(s) > 0$, then $s$ lies on the critical line $\sigma = 1/2$ if and only if
\begin{equation}
  \ReOp\bigl( \overline{S_{N-1}(s)}\, w_N(s) \bigr) \le 0 \quad \text{for all } N \ge 1,
  \label{eq:greedy_disc}
\end{equation}
where $S_N(s) = \sum_{n=1}^N (-1)^{n-1} n^{-s}$ are the partial sums and $w_N(s) = N^{-s}$.  This condition has a direct translation into the language of the iota function.

For the continuous version, define the accumulated projected sum
\[
  \Sigma(\tau) = \int_1^\tau \e^{\iu\pi(u-1)} u^{-s}\,du.
\]
Then the continuous analogue of \eqref{eq:greedy_disc} is
\begin{equation}
  \ReOp\Bigl( \overline{\Sigma(\tau)}\, \tau^{-s} \e^{\iu\pi(\tau-1)} \Bigr) \le 0
  \quad \text{for all } \tau \ge 1.
  \label{eq:greedy_cont}
\end{equation}

\begin{proposition}
  A non‑trivial zero $s = \sigma + \iu t$ of $\eta(s)$ satisfies the Riemann Hypothesis (i.e., $\sigma = 1/2$) if and only if the continuous iota curve $\iota_c(s)$ satisfies the geometric inequality \eqref{eq:greedy_cont}.
\end{proposition}

Geometrically, the inequality states that at every instant $\tau$, the tangent vector to the curve $\Sigma(\tau)$ in the complex plane never has a positive component along the direction of the instantaneous vector $\tau^{-s}\e^{\iu\pi(\tau-1)}$ (which is precisely the horizontal projection of $\iota_c(\tau)$).  In other words, the spiral $\Sigma(\tau)$ always ``turns inward'' relative to its current position.

\subsection{Connection to the Functional Equation}
The functional equation of $\zeta(s)$,
\[
  \zeta(s) = 2^s \pi^{s-1} \sin\Bigl(\frac{\pi s}{2}\Bigr) \Gamma(1-s) \zeta(1-s),
\]
induces a symmetry on the iota curve when $s$ is replaced by $1-s$.  For $s$ on the critical line $\sigma = 1/2$, the curve exhibits a self‑similarity that reflects the approximate functional equation of the zeta function.  This suggests that geometric invariants of the iota spiral (such as total curvature or a renormalised torsion) might characterise the critical line.

\subsection{Unification of Known Approaches to the Riemann Hypothesis}\label{subsec:unification}

The iota function provides a single geometric framework that encapsulates several classical equivalent formulations of the Riemann Hypothesis.  By translating abstract conditions on function spaces or variational principles into concrete properties of the iota spiral, the function unifies the following approaches:

\begin{enumerate}
  \item \textbf{Nyman--Beurling Criterion.}  The Nyman--Beurling theorem states that the Riemann Hypothesis is true if and only if the characteristic function of the interval $(0,1]$ belongs to the closed linear span in $L^2(0,\infty)$ of the set of dilations $\{ \rho_\alpha(x) = \{\alpha/x\} \mid 0 < \alpha \le 1 \}$, where $\{\cdot\}$ denotes the fractional part.  In the language of the iota function, the relevant dilations correspond precisely to the horizontal projections of the continuous spiral $\iota_c(\tau)$ at scaled parameters.  The approximation of $\chi_{(0,1]}$ by finite linear combinations of these dilations translates into a statement about how well the projected curve $\Sigma(\tau)$ can approximate a step function—a problem that is visually and analytically accessible once the curve is realised geometrically.

  \item \textbf{B\'aez-Duarte Criterion.}  B\'aez-Duarte strengthened the Nyman--Beurling result by replacing the continuous dilations with a discrete set.  His criterion asserts that RH holds if and only if the characteristic function of $(0,1]$ lies in the closed linear span of the sequence $\{ \rho_n(x) = \{1/(nx)\} \}_{n=1}^\infty$.  The discrete iota points $P_n$ are the exact geometric counterpart of these dilations: the $x$ and $y$ coordinates of $P_n$ are the real and imaginary parts of the $n$-th term of the eta series, which are intimately related to the fractional‑part functions via the Mellin transform.  Thus the B\'aez-Duarte condition becomes a statement about the density of the sequence $\{P_n\}$ in a suitable function space constructed from the spiral.

  \item \textbf{The Greedy Condition.}  As already detailed in \S\ref{subsec:greedy}, the greedy condition \eqref{eq:greedy_disc} is an exact equivalent of RH for the eta function.  The iota function translates this sequential inequality into a dynamic, pointwise condition \eqref{eq:greedy_cont} on the tangent vector of the projected spiral.  This geometric reformulation reveals that the Riemann Hypothesis is equivalent to a \emph{stability property} of the continuous iota curve: the accumulated displacement never points ``outward'' from the origin relative to the instantaneous term.  Such a condition is reminiscent of Lyapunov stability criteria in dynamical systems.

  \item \textbf{Bombieri's Variational Principle.}  Bombieri (following work of Weil and Guinand) showed that the Riemann Hypothesis is equivalent to the positivity of a certain quadratic functional—the Weil explicit formula.  This functional can be expressed as a variational problem over test functions.  Through the iota lens, the test functions correspond to smooth deformations of the iota spiral, and the quadratic functional becomes an energy integral along the curve.  Minimising this functional is equivalent to finding the spiral that best satisfies the constraints imposed by the prime numbers.  The critical line $\sigma = 1/2$ emerges as the unique parameter for which the spiral achieves a critical point of this geometric action.
\end{enumerate}

In each case, the iota function serves as a \emph{universal translator}, converting the distinct languages of functional analysis, discrete approximation, sequential inequalities, and variational calculus into a common geometric vocabulary.  This unification does not, by itself, prove the Riemann Hypothesis; it does, however, reveal deep structural connections between the various approaches and provides a coherent picture in which progress in one area can be directly interpreted in the others.

\section{Applications and Impact in Mathematical Fields}

\subsection{Complex Analysis and Analytic Number Theory}
The iota function provides a concrete geometric model for the behaviour of Dirichlet series.  It translates subtle problems of exponential sums and oscillatory integrals into questions about the shape of a single curve.  This may offer new perspectives on:
\begin{itemize}
  \item \textbf{Exponential sum estimates}: The phase $\phi(\tau) = \pi(\tau-1) - t\ln\tau$ encodes the distribution of $t\ln n \bmod 2\pi$.  Bounding the partial sums of the eta series is equivalent to controlling the net displacement of the projected spiral.
  \item \textbf{Zero‑density theorems}: The injectivity of $\iota$ means that distinct zeros produce distinct spirals.  Statistical properties of these spirals could be linked to the distribution of zeros.
  \item \textbf{Universality}: Voronin's universality theorem for $\zeta(s)$ might be rephrased as a statement about the density of the iota curves in a suitable function space.
\end{itemize}

\subsection{Dynamical Systems}
The continuous iota curve $\Sigma(\tau)$ can be viewed as the trajectory of a particle in the complex plane subject to a non‑autonomous force field.  The greedy condition \eqref{eq:greedy_cont} resembles a Lyapunov stability criterion: the origin is a stable attractor for the dynamical system defined by the differential equation
\[
  \frac{d\Sigma}{d\tau} = \e^{\iu\pi(\tau-1)} \tau^{-s}.
\]
This opens the possibility of applying techniques from stability theory and control theory to the study of RH.

\subsection{Differential Geometry and Special Functions}
The iota curve is a parametric plot of the \emph{incomplete gamma function} (after a change of variables).  Indeed,
\[
  \int_1^\tau \e^{\iu\pi u} u^{-s}\,du = (-\iu\pi)^{s-1} \Bigl( \Gamma(1-s, -\iu\pi) - \Gamma(1-s, -\iu\pi\tau) \Bigr).
\]
Thus the iota function links the Riemann zeta function to the monodromy and asymptotics of the incomplete gamma function.  This connection is already implicit in the Riemann--Siegel formula, but the geometric interpretation makes it more tangible.

\subsection{Numerical and Visual Exploration}
The smoothness of $\iota_c$ makes it amenable to numerical integration and visualisation.  One can plot the curve for various $s$, observe the violation of the greedy condition off the critical line, and potentially discover new invariants.  Such visual experiments have heuristic value and may inspire new conjectures.

\subsection{Potential Impact on the Riemann Hypothesis}
While the iota function does not, by itself, resolve the Riemann Hypothesis, it provides a rigorous and elegant framework for restating the problem.  Any geometric proof of RH using this framework would still have to confront the deep arithmetic nature of the phase $t\ln\tau$.  Nevertheless, the iota function unifies several known approaches (Nyman--Beurling, B\'aez-Duarte, Bombieri) under a single geometric umbrella, which may facilitate cross‑fertilisation of ideas.

\section{Conclusion}

We have introduced the iota function—both in discrete and continuous forms—as an invertible geometric encoding of the Dirichlet eta series.  The function maps a complex number $s$ to a three‑dimensional spiral whose horizontal projection directly corresponds to the terms of the series.  Injectivity guarantees that no information is lost, allowing statements about $\eta(s)$ to be translated into purely geometric conditions.

The continuous iota curve reveals a deep connection with the incomplete gamma function and provides a new perspective on the Riemann Hypothesis through the greedy condition and the unification of several classical approaches.  Beyond number theory, the iota function may find applications in complex analysis, dynamical systems, and the visualisation of special functions.  While it does not constitute a proof of RH, it offers a promising geometric language that may inspire novel approaches to this long‑standing problem.

\begin{thebibliography}{9}
\bibitem{BaezDuarte}
L.~B\'aez-Duarte, \emph{A class of invariant unitary operators}, Adv. Math. \textbf{144} (1999), 1--12.

\bibitem{Bombieri}
E.~Bombieri, \emph{The Riemann Hypothesis}, in: \emph{The Millennium Prize Problems}, Clay Math. Inst., 2006.

\bibitem{Nyman}
B.~Nyman, \emph{On the One‑Dimensional Translation Group and Semi‑Group in Certain Function Spaces}, Thesis, Uppsala, 1950.

\bibitem{Titchmarsh}
E.~C.~Titchmarsh, \emph{The Theory of the Riemann Zeta‑function}, 2nd ed., Oxford Univ. Press, 1986.

\bibitem{Beurling}
A.~Beurling, \emph{A closure problem related to the Riemann zeta‑function}, Proc. Nat. Acad. Sci. U.S.A. \textbf{41} (1955), 312--314.

\end{thebibliography}

\end{document}